\par\noindent\begin{minipage}{\linewidth}
\tagstructbegin{tag=Figure,alttext={同用两张卡的四种分工。数据并行各有完整模型，分别处理样本并同步梯度；状态分片各保存半份状态并按需通信；流水线卡一保存第一二层、卡二保存第三四层并传激活；张量并行各保存同一矩阵一半输出列并拼接结果。}}
\begin{center}
\begin{tikzpicture}[x=1mm,y=1mm,>=Latex,
 device/.style={draw=linegray,rounded corners=2pt,fill=paper,align=center,font=\footnotesize,text width=39mm,minimum height=12mm}]
\node[font=\small\bfseries] at (45,10) {卡 1};
\node[font=\small\bfseries] at (111,10) {卡 2};
\node[font=\footnotesize,anchor=east] at (19,0) {数据并行};
\node[device] (d1) at (45,0) {完整模型\\样本前半批};
\node[device] (d2) at (111,0) {完整模型\\样本后半批};
\draw[<->,thick] (d1)--node[above,font=\scriptsize]{梯度}(d2);
\node[font=\footnotesize,anchor=east] at (19,-19) {状态分片};
\node[device] (s1) at (45,-19) {状态分片 1};
\node[device] (s2) at (111,-19) {状态分片 2};
\draw[<->,thick] (s1)--node[above,font=\scriptsize]{按需取用}(s2);
\node[font=\footnotesize,anchor=east] at (19,-38) {流水线并行};
\node[device] (p1) at (45,-38) {第 1、2 层};
\node[device] (p2) at (111,-38) {第 3、4 层};
\draw[->,thick] (p1)--node[above,font=\scriptsize]{激活}(p2);
\node[font=\footnotesize,anchor=east] at (19,-57) {张量并行};
\node[device] (t1) at (45,-57) {同一层权重第 1、2 列};
\node[device] (t2) at (111,-57) {同一层权重第 3、4 列};
\draw[<->,thick] (t1)--node[above,font=\scriptsize]{拼接输出}(t2);
\end{tikzpicture}
\captionof{figure}{先问拆什么：数据、训练状态、完整层，还是一层中的矩阵}
\label{fig:distributed-two-devices}
\end{center}
\tagstructend
\end{minipage}\par
